% -------------------------------------------------------------------------------------
\section{领头阶共线拉氏量与共线 Wilson 线}\label{sec:scet-lagrangian}

本节先在树图层次把 QCD 夸克场分解为大、小两个分量并积掉小分量——这一步没有任何近似——再按\cref{prop:scet-scalings}去掉被幂次压低的超软耦合，得到领头阶共线拉氏量；然后说明共线 Wilson 线 $W_n$ 如何从积掉离壳的硬线中产生。令 $\xi_n\equiv P_n\psi$、$\phi_{\bar n}\equiv P_{\bar n}\psi$，则 $\sla n\,\xi_n=0$，$\bar\xi_n=\bar\psi P_{\bar n}$；对 $n$ 共线夸克，$\phi_{\bar n}$ 是被 $\lambda$ 压低的分量。

\begin{proposition}[领头阶共线夸克拉氏量 \tagR]\label{prop:scet-lagrangian}
（$R_{\mathrm{PT}}$：在重整化微扰论中对红外正规化的部分子态严格。）（树图层次的匹配。）
(a) 在经典作用量 $\int\bar\psi\,\ii\sla D\,\psi$（\cref{def:prelim-qcd}）中解出 $\phi_{\bar n}$ 的运动方程并代回，精确地得到
\begin{equation}\label{eq:scet-Lxixi-exact}
\mathcal L_{\xi\xi}=\bar\xi_n\Big[\ii n\cdot D+\ii\sla D_\perp\,\frac{1}{\ii\bar n\cdot D}\,\ii\sla D_\perp\Big]\frac{\bar{\sla n}}{2}\,\xi_n .
\end{equation}
(b) 令 $A=A_n+A_{us}$，超软场按\cref{eq:scet-multipole}多极展开并只保留领头幂次，得
\begin{equation}\label{eq:scet-L0xi}
\begin{gathered}
\mathcal L^{(0)}_{\xi\xi}=\bar\xi_n\Big[\ii n\cdot D+\ii\sla D_{\perp c}\,\frac{1}{\ii\bar n\cdot D_c}\,\ii\sla D_{\perp c}\Big]\frac{\bar{\sla n}}{2}\,\xi_n,\\
\ii n\cdot D=\ii n\cdot\partial+g\,n\cdot A_n+g\,n\cdot A_{us}(x_-),\\
\ii\bar n\cdot D_c=\ii\bar n\cdot\partial+g\,\bar n\cdot A_n,\qquad \ii D^\mu_{\perp c}=\ii\partial^\mu_\perp+gA^\mu_{n\perp}.
\end{gathered}
\end{equation}
% VF-SCET-3
(c) Feynman 规则：共线夸克传播子 $S_\xi(p)=\ii\frac{\sla n}{2}\frac{\bar n\cdot p}{p^2+\ii0}$；超软胶子顶角 $\ii g\,t^a\,n^\mu\frac{\bar{\sla n}}{2}$；共线胶子顶角（夸克 $p\to p'$）
\begin{equation}\label{eq:scet-collvertex}
V^\mu=\ii g\,t^a\Big[n^\mu+\frac{\gamma_\perp^\mu\sla p_\perp}{\bar n\cdot p}+\frac{\sla p'_\perp\gamma_\perp^\mu}{\bar n\cdot p'}-\frac{\sla p'_\perp\sla p_\perp}{\bar n\cdot p\,\bar n\cdot p'}\,\bar n^\mu\Big]\frac{\bar{\sla n}}{2},\qquad
\gamma_\perp^\mu=\gamma^\mu-n^\mu\frac{\bar{\sla n}}{2}-\bar n^\mu\frac{\sla n}{2};
\end{equation}
% VF-SCET-4
并且对任意离壳动量精确成立 $P_n\,S(p')\,(\ii g\gamma^\mu t^a)\,S(p)\,P_{\bar n}=S_\xi(p')\,V^\mu\,S_\xi(p)$，$S(p)=\frac{\ii\sla p}{p^2+\ii0}$。（本命题只涉及部分子场与部分子振幅。）
\uses{prop:scet-scalings, def:prelim-qcd}
\end{proposition}

\begin{proof}
(a) 把 $\sla D=\frac{\sla n}{2}\bar n\cdot D+\frac{\bar{\sla n}}{2}n\cdot D+\sla D_\perp$ 代入，投影算符的性质给出 $P_{\bar n}\sla DP_n=n\cdot D\frac{\bar{\sla n}}{2}P_n$、$P_n\sla DP_{\bar n}=\bar n\cdot D\frac{\sla n}{2}P_{\bar n}$、$P_{\bar n}\sla DP_{\bar n}=\sla D_\perp P_{\bar n}$、$P_n\sla DP_n=\sla D_\perp P_n$；由 $\bar\phi_{\bar n}=\bar\psi P_n$，
\begin{equation}\label{eq:scet-psiDpsi}
\bar\psi\,\ii\sla D\,\psi=\bar\xi_n\frac{\bar{\sla n}}{2}\,\ii n\cdot D\,\xi_n+\bar\xi_n\,\ii\sla D_\perp\phi_{\bar n}
+\bar\phi_{\bar n}\,\ii\sla D_\perp\xi_n+\bar\phi_{\bar n}\frac{\sla n}{2}\,\ii\bar n\cdot D\,\phi_{\bar n}.
\end{equation}
对 $\bar\phi_{\bar n}$ 变分得 $\frac{\sla n}{2}\ii\bar n\cdot D\,\phi_{\bar n}+\ii\sla D_\perp\xi_n=0$，左乘 $\frac{\bar{\sla n}}{2}$ 并用 $\frac{\bar{\sla n}\sla n}{4}\phi_{\bar n}=\phi_{\bar n}$ 得 $\phi_{\bar n}=-\frac{1}{\ii\bar n\cdot D}\frac{\bar{\sla n}}{2}\,\ii\sla D_\perp\xi_n$。代回后含 $\bar\phi_{\bar n}$ 的两项按运动方程为零；把余项中的 $\bar{\sla n}$ 越过 $\sla D_\perp$ 移到最右（多一个负号）即得\cref{eq:scet-Lxixi-exact}。这一步没有近似（路径积分中对 $\phi_{\bar n}$ 积分只另给出一个与 $\xi_n$ 无关的行列式）。
(b) 作用在共线场上时 $\ii n\cdot\partial$、$g\,n\cdot A_n$、$g\,n\cdot A_{us}$ 都是 $\order{\lambda^2}$；$\ii\bar n\cdot\partial$、$g\,\bar n\cdot A_n\sim1$ 而 $g\,\bar n\cdot A_{us}\sim\lambda^2$；$\ii\partial_\perp$、$gA_{n\perp}\sim\lambda$ 而 $gA_{us\perp}\sim\lambda^2$（\cref{prop:scet-scalings}）。去掉被压低的项即得\cref{eq:scet-L0xi}。
(c) 在 $P_n$ 子空间上自由部分等于 $n\cdot p+\sla p_\perp\sla p_\perp/\bar n\cdot p=p^2/\bar n\cdot p$，而 $S_\xi\,\frac{p^2}{\bar n\cdot p}\frac{\bar{\sla n}}{2}=\ii P_n$，故传播子如上（即\cref{eq:scet-xiprop}）。超软顶角来自 $\bar\xi_n\,g\,n\cdot A_{us}\frac{\bar{\sla n}}2\xi_n$。共线顶角是 $A_{us}=0$ 时的 $\order{g}$ 项：$n^\mu$ 来自 $\ii n\cdot D$，中间两项来自 $\ii\sla D_{\perp c}$ 中的 $gA_{n\perp}$，末项来自 $(\bar n\cdot p+g\bar n\cdot A_n)^{-1}=\frac{1}{\bar n\cdot p}-\frac{1}{\bar n\cdot p'}g\,\bar n\cdot A_n\frac{1}{\bar n\cdot p}+\order{g^2}$。由 (a)，消去 $\phi_{\bar n}$ 后的树图生成泛函等于 QCD 树图生成泛函对 $\xi_n$ 源的限制，而 $A_{us}=0$ 时 (b) 没有丢掉任何项，故末一恒等式精确成立（对随机离壳动量逐元素核对，相对差 $<10^{-15}$）。% w18_s1 (D)
\end{proof}

\paragraph{共线 Wilson 线。}
由于 $g\,\bar n\cdot A_n\sim\lambda^0$，定义
\begin{equation}\label{eq:scet-Wn}
W_n(x)\equiv W[x,x-\infty\bar n]=\mathrm P\exp\Big(\ii g\int_{-\infty}^0\dd s\;\bar n\cdot A_n(x+s\bar n)\Big)
=\sum_{\rm perms}\exp\Big[-\frac{g}{\bar{\mathcal P}}\,\bar n\cdot A_n\Big]
\end{equation}
% VF-SCET-5
（后一式是标签动量表述中的形式\cite{BauerEtAl2001,BauerPirjolStewart2001}，$\bar{\mathcal P}$ 作用于其右边的场，对胶子标签的排列求和）。由\cref{prop:prelim-wilson-covariance}，$W_n^\dagger W_n=1$，$\bar n\cdot D_cW_n=0$，故 $\frac{1}{\ii\bar n\cdot D_c}=W_n\frac{1}{\ii\bar n\cdot\partial}W_n^\dagger$；共线规范变换（$U_c(x-\infty\bar n)=1$）下 $W_n\to U_c(x)W_n$，故 $\chi_n\equiv W_n^\dagger\xi_n$ 共线规范不变。被吸收的胶子（场 $\propto\epsilon\,\ee^{-\ii k\cdot y}$）给出 $\ii g\,\bar n\cdot\epsilon\,t^a\int_{-\infty}^0\dd s\,\ee^{-\ii s\bar n\cdot k+0s}=-g\,t^a\bar n\cdot\epsilon/(\bar n\cdot k+\ii0)$，出射胶子给出 $+g\,t^a\bar n\cdot\epsilon^*/(\bar n\cdot k-\ii0)$，即\cref{prop:prelim-wilson-rules}中指向过去的线的规则（$n\to\bar n$）。

%%% BEGIN fig:scet-wilson
%%% (replaced verbatim from figblocks.tex at assembly)
%%% END fig:scet-wilson

\begin{proposition}[共线 Wilson 线 \tagR]\label{prop:scet-wilson}
（$R_{\mathrm{PT}}$：在重整化微扰论中对红外正规化的部分子态严格。）（树图层次。）设 QCD 矢量流 $J^\mu=\bar\psi\gamma^\mu\psi$ 产生一个 $n$ 共线夸克、一个 $\bar n$ 共线反夸克及任意个 $n$ 共线与 $\bar n$ 共线胶子。积掉虚度 $\sim Q^2$ 的离壳线后，其树图矩阵元在 $\lambda$ 的领头幂次等于
\begin{equation}\label{eq:scet-Jtree}
J^\mu_{\rm SCET}=\bar\xi_n\,W_n\,\gamma_\perp^\mu\,W_{\bar n}^\dagger\,\xi_{\bar n}=\bar\chi_n\gamma_\perp^\mu\chi_{\bar n}
\end{equation}
用 $\mathcal L^{(0)}$ 算出的矩阵元，$W_n$ 由\cref{eq:scet-Wn}给出，单胶子规则为 $-g\,\bar n^\mu t^a/(\bar n\cdot k+\ii0)$。与 BFPS、BPS 相比不需要改变 $g$ 的符号。（本命题只涉及部分子振幅。）
\uses{prop:scet-scalings, prop:scet-lagrangian, def:prelim-wilson, prop:prelim-wilson-covariance, prop:prelim-wilson-rules, lem:prelim-eikonal-identity}
\end{proposition}

\begin{proof}
设胶子都是 $n$ 共线的（$\bar n$ 共线胶子同理给出 $W_{\bar n}^\dagger$）。外胶子的极化矢量取\cref{prop:scet-scalings}的标度 $(\bar n\cdot\epsilon,n\cdot\epsilon,\epsilon_\perp)\propto(1,\lambda^2,\lambda)$：这总可由 $\epsilon\to\epsilon+c\,k$ 做到，而 QCD 振幅（Ward 恒等式）与 SCET 振幅（共线规范不变性）都不变。\emph{分解。}任一树图中，从流顶点到外反夸克的线携带 $p_2+K$，$K$ 为若干 $n$ 共线动量之和，$(p_2+K)^2=n\cdot p_2\,\bar n\cdot K[1+\order{\lambda^2}]\sim Q^2$：这些是要积掉的离壳线。接在其上的每个胶子是一棵只含 $n$ 共线线的子树之根；子树与 $n$ 共线夸克线属于共线扇区，由\cref{prop:scet-lagrangian}精确描述。\emph{领头幂次。}设反夸克线上自外反夸克起依次接了 $r$ 棵子树（流 $\mathcal J_j$，色指标 $a_j$，$K_j$ 为前 $j$ 棵的总动量），反夸克线给出
\[
\Gamma^\mu\,\frac{\ii(-\sla p_2-\sla K_r)}{(p_2+K_r)^2+\ii0}\,(\ii g\,\sla{\mathcal J}_r t^{a_r})\cdots\frac{\ii(-\sla p_2-\sla K_1)}{(p_2+K_1)^2+\ii0}\,(\ii g\,\sla{\mathcal J}_1t^{a_1})\,v(p_2).
\]
领头幂次下 $\sla{\mathcal J}_j\to\frac{\bar n\cdot\mathcal J_j}{2}\sla n$（$n\cdot\mathcal J_j\sim\lambda^2$，$\mathcal J_{j\perp}\sim\lambda$），而 $(-\sla p_2-\sla K_j)\sla n=-\frac{n\cdot p_2}{2}\bar{\sla n}\sla n+\order{\lambda}$（$\sla n^2=0$），故每对“传播子 × 顶角”化为 $\frac{g\,\bar n\cdot\mathcal J_j}{4(\bar n\cdot K_j+\ii0)}\bar{\sla n}\sla n\,t^{a_j}$；由 $\bar{\sla n}v(p_2)=\order{\lambda}$ 得 $\bar{\sla n}\sla n\,v=4v+\order{\lambda}$，整条线化为 $\Gamma^\mu v(p_2)\prod_{j=r}^{1}[g\,t^{a_j}\bar n\cdot\mathcal J_j/(\bar n\cdot K_j+\ii0)]$（$j=r$ 在最左），且 $\bar u_n\gamma^\mu v_{\bar n}=\bar u_n\gamma_\perp^\mu v_{\bar n}$。\emph{求和。}对子树的根在线上的 $r!$ 种次序求和。另一方面，\cref{eq:scet-Wn}对 $r$ 个出射胶子在 $s_1<\cdots<s_r\le0$ 上的嵌套积分逐层可积，$\int_{-\infty}^{s_{j+1}}\dd s_j\,\ee^{\ii s_j\bar n\cdot K_j}=\ee^{\ii s_{j+1}\bar n\cdot K_j}/(\ii\bar n\cdot K_j)$，给出 $(\ii g)^r\prod_j(\ii\bar n\cdot K_j)^{-1}=g^r/\prod_j\bar n\cdot K_j$，色矩阵按 $t^{a_r}\cdots t^{a_1}$ 排列，恰为上式在同一次序下的值；对次序求和即路径有序指数的展开（\cref{prop:prelim-wilson-rules}；交换群极限即\cref{lem:prelim-eikonal-identity}）。QCD 离壳线给出 $\bar n\cdot K_j+\ii0$，$W_n$ 给出 $-\ii0$，但共线胶子有 $\bar n\cdot K_j\sim Q\neq0$；$W_n$ 取自 $-\infty$ 是约定：方向依赖只来自 $\bar n\cdot k\to0$ 的区域，在 SCET 中被共线 Wilson 线图的软与 Glauber 零箱减除抵消\cite[第~4.2~节]{RothsteinStewart2016}。于是每个 QCD 树图等于子树接在 $W_n$ 的 $r$ 胶子顶点上的 SCET 图，即\cref{eq:scet-Jtree}。推导只用 $D=\partial-\ii gA$ 与 $W=\mathrm P\exp(+\ii g\int\dd z\cdot A)$，与 BFPS、BPS 的约定相同。
\end{proof}


\begin{example}[$W_n$ 的树图匹配]\label{ex:scet-wn-tree}
取 $\gamma^*\to q(p_1)\bar q(p_2)g(k)$，$k\parallel p_1\parallel n$（\cref{fig:scet-wilson}(a)）。胶子接在夸克上的图由共线顶角描述（\cref{prop:scet-lagrangian}(c)）；接在反夸克上的图 $\bar u(p_1)\gamma^\mu\frac{\ii(-\sla p_2-\sla k)}{(p_2+k)^2+\ii0}(\ii g\,t^a\sla\epsilon^*)v(p_2)$ 中 $(p_2+k)^2\simeq n\cdot p_2\,\bar n\cdot k\sim Q^2$，在领头幂次等于 $W_n$ 的单胶子项 $g\,t^a\frac{\bar n\cdot\epsilon^*}{\bar n\cdot k}\bar u(p_1)\gamma_\perp^\mu v(p_2)$。对流的横向分量与一个 $\bar n\cdot\epsilon^*=\order{1}$ 的物理极化，在随机运动学点上二者之差的相对大小在 $\lambda=10^{-1},10^{-2},10^{-3}$ 时为 $4.6\times10^{-2}$、$4.7\times10^{-3}$、$4.8\times10^{-4}$；完整的 QCD 与 SCET 振幅之差为 $9.1\times10^{-3}$、$8.5\times10^{-5}$、$8.5\times10^{-7}$。% w18_s2 (B)
取 $\epsilon^*\to k$：QCD 两图之和为零；SCET 中共线顶角项 $-g\,t^a\bar u\gamma_\perp^\mu v$ 与 $W_n$ 项 $+g\,t^a\bar u\gamma_\perp^\mu v$ 精确相消——$W_n$ 恢复了共线规范不变性。
\uses{prop:scet-wilson}
\end{example}
